\documentclass[aspectratio=169]{beamer}
\usetheme{Madrid}
\usecolortheme{default}
\usefonttheme{professionalfonts}

\usepackage[utf8]{inputenc}
\usepackage[T1]{fontenc}
\usepackage[spanish,es-nodecimaldot]{babel}
\usepackage{lmodern}
\usepackage{ragged2e}
\usepackage{booktabs}
\usepackage{enumitem}
\usepackage{hyperref}

\setbeamertemplate{navigation symbols}{}
\setbeamertemplate{footline}[frame number]
\setbeamertemplate{bibliography item}[text]

\title{Riesgos en la Elaboración del Plan de Tesis}
\subtitle{Trabajo de Investigación I\\MIA 3er Ciclo A 2026-1\\Grupo 9B}
\author{Luis Koc \\ Herbert Meléndez}
\date{\today}

\newcommand{\smallrefs}[1]{{\footnotesize #1}}

\begin{document}

\begin{frame}
  \titlepage
\end{frame}

\begin{frame}{Tema propuesto}
\textbf{{\large Efecto de la heterogeneidad térmica del entorno sobre la
temperatura del conductor en cables subterráneos: modelado numérico 2D y
análisis de la brecha normativa}}

\vspace{0.5em}
\small
\textit{Variable X:} conductividad térmica del suelo/relleno $k(x,y)$ variable
\quad$\longrightarrow$\quad
\textit{Variable Y:} temperatura del conductor $T_{\text{cond}}$ / ampacidad estimada

\vspace{0.6em}
\justifying
Los cables eléctricos subterráneos son infraestructura crítica cuya capacidad
de conducción depende del comportamiento térmico del entorno de instalación.
Los métodos vigentes (IEC~60287, Neher-McGrath) asumen conductividad térmica
uniforme y condiciones estacionarias; cuando el entorno real presenta
heterogeneidad —por variación de humedad, compactación, tipo de relleno o
temperatura estacional del suelo— esos modelos sobrestiman o subestiman la
temperatura real del conductor, generando una \textbf{brecha entre la
capacidad nominal y la capacidad efectiva} del sistema.

\vspace{0.3em}
\justifying
La investigación busca caracterizar esa brecha mediante un modelo numérico 2D
de la ecuación de calor $\nabla\cdot(k(x,y)\,\nabla T)+Q=0$ con $k$
espacialmente variable, resolver también el \textbf{problema inverso}
(estimar $k$ del suelo a partir de mediciones de temperatura) y desarrollar
una herramienta adaptable, de código abierto, validada con casos de referencia
publicados.
\end{frame}

\section{Riesgos de ejecución}

\begin{frame}{Perfil de los tesistas y viabilidad de ejecución}
\justifying
El perfil de ambos tesistas muestra una combinación de fortalezas para una tesis aplicada: experiencia profesional extensa en energía, conocimientos del idioma inglés, capacidad de programación y familiaridad con herramientas digitales.

\vspace{0.5em}
No obstante, precisamente por tratarse de perfiles ejecutivo-técnicos, la realización de la tesis enfrenta riesgos específicos que deben ser reconocidos desde el inicio como parte de la planificación académica.
\end{frame}

\begin{frame}{Riesgos principales de ejecución de la tesis}
\justifying
\begin{itemize}[leftmargin=*]
    \item \textbf{Restricción de tiempo y continuidad:} ambos tesistas ocupan cargos de responsabilidad a tiempo completo que compiten directamente con los bloques de trabajo académico sostenido requeridos para avanzar en la tesis.
    \item \textbf{Fragmentación atencional:} los múltiples frentes operativos de cada tesista reducen la concentración necesaria para desarrollar una investigación rigurosa y coherente.
    \item \textbf{Ampliación excesiva del alcance:} la experiencia técnica amplia de ambos puede llevar a formular un problema demasiado ambicioso o multidimensional para los plazos de una maestría.
    \item \textbf{Sesgo hacia la solución prematura:} el dominio técnico compartido genera tendencia a diseñar la herramienta o el modelo antes de haber delimitado con precisión el problema de investigación.
    \item \textbf{Riesgo de confidencialidad y acceso a datos:} aspectos del tema pueden depender de datos de instalación o de operación de empresas que no los hacen públicos ni permiten su uso en publicaciones académicas.
    \item \textbf{Ajuste al registro académico:} la escritura técnico-profesional no siempre satisface las exigencias de argumentación, citación y estructura propias de una tesis de posgrado.
\end{itemize}
\end{frame}

\begin{frame}{Riesgos adicionales vinculados al tema específico}
\justifying
\begin{itemize}[leftmargin=*]
    \item \textbf{Disponibilidad de datos de campo:} los parámetros térmicos del suelo (resistividad, humedad, temperatura) en instalaciones reales rara vez están publicados con el detalle requerido para validar modelos, lo que puede limitar la profundidad del análisis.
    \item \textbf{Complejidad interdisciplinaria:} el tema exige dominio simultáneo de transferencia de calor, geotecnia, normativa eléctrica e inteligencia artificial, elevando el riesgo de tratamiento superficial en alguna de estas áreas.
    \item \textbf{Integración conceptual de los dos ejes del problema:} articular el comportamiento térmico variable con la observabilidad de los cuellos de botella requiere una delimitación cuidadosa para evitar una tesis con dos problemas en lugar de uno.
    \item \textbf{Dependencia del estado del arte emergente:} parte de la literatura relevante es reciente (2024--2025); esto representa una oportunidad, pero también un riesgo si las referencias aún no están suficientemente consolidadas.
\end{itemize}
\end{frame}

\begin{frame}{Factores mitigantes y controles recomendados}
\justifying
El perfil de ambos tesistas contiene factores mitigantes relevantes: disciplina profesional, madurez personal, base analítica sólida y capacidad para revisar literatura internacional. Adicionalmente, el trabajo en equipo permite distribuir la carga de revisión bibliográfica y escritura.

\vspace{0.5em}
Se recomienda:
\begin{itemize}[leftmargin=*]
    \item reservar bloques horarios fijos e inamovibles para la tesis,
    \item mantener una sola pregunta central de investigación que articule ambos ejes del problema,
    \item separar la etapa de comprensión del problema de la etapa metodológica,
    \item elegir un tema sostenible con datos públicos o anonimizables, priorizando fuentes de acceso abierto,
    \item trabajar con hitos mensuales breves y acumulativos, y
    \item delimitar desde el inicio qué componente de IA responde a cuál aspecto del problema.
\end{itemize}

\vspace{0.5em}
\textbf{Conclusión:} la tesis es viable para este equipo si se gestiona con rigor el tiempo disponible de cada tesista y, sobre todo, la delimitación del alcance del problema.
\end{frame}

\begin{frame}
\centering
\Large Gracias
\end{frame}

\end{document}
